\section*{Spectral Efficiency}
\begin{equation}
\eta \defeq \frac {R} {\BW}
\label{eq:spectral_efficiency}
\end{equation}

\section*{QAM Signal}

\paragraph{QAM modulated signal}
\begin{equation}
x_{\textrm{QAM}}(t) = x_i(t) \cos(\aw_c t) + x_q(t) \sin(\aw_c t)
\label{eq:qam}
\end{equation}

\paragraph{Alternative convention (subtraction)}
\begin{equation}
x_{\textrm{QAM}}(t) = x_i(t) \cos(\aw_c t) - x_q(t) \sin(\aw_c t)
\label{eq:qam2}
\end{equation}

\paragraph{Recovering the in-phase component}
\begin{eqnarray}
x_{\textrm{QAM}}(t) \cos(\aw_c t) & = & x_i(t) \cos^2(\aw_c t) - x_q(t) \sin(\aw_c t)  \cos(\aw_c t) \nonumber \\
& = & \frac{1}{2} \left[ x_i(t)(1+\cos(2\aw_c t)) - x_q(t) \sin(2\aw_c t) \right] \nonumber \\
& = & \frac{x_i(t)}{2} + \frac{1}{2} \left[ x_i(t)\cos(2\aw_c t) - x_q(t) \sin(2\aw_c t)\right]
\label{eq:qam_recovery}
\end{eqnarray}

\paragraph{Recovery via integration (correlation)}
\begin{eqnarray*}
\int_{-\infty}^\infty x_{\textrm{QAM}}(t) \cos(\aw_c t) dt & = & \int_{-\infty}^\infty \left[ x_i(t) \cos^2(\aw_c t) - x_q(t) \sin(\aw_c t)  \cos(\aw_c t) \right] dt\\
& = & \frac{1}{2} \int_{-\infty}^\infty x_{i}(t) dt
\end{eqnarray*}

\section*{Complex Envelope}

\paragraph{Definition (matching the ``+'' convention)}
\[
x_{\textrm{ce}}(t) = x_i(t) - j x_q(t)
\]
\begin{align}
x_{\textrm{QAM}}(t) = & \real{ x_{\textrm{ce}}(t) e^{j \aw_c t} } \notag \\
 = & \real{ [x_i(t)- jx_q(t)] [\cos(\aw_c t)+j\sin(\aw_c t)]} \notag \\
 = & x_i(t) \cos(\aw_c t) + x_q(t) \sin(\aw_c t)
\label{eq:qam_complex}
\end{align}

\paragraph{Definition (matching the ``$-$'' convention)}
\begin{equation}
x_{\textrm{ce}}(t) = x_i(t) + j x_q(t)
\label{eq:complex_envelope}
\end{equation}

\paragraph{Spectrum of the QAM signal}
\begin{align}
X_{\textrm{QAM}}(\aw) = & \calF \{ x_{\textrm{QAM}}(t) \}  \notag \\
 = & \calF \{ \real{ x_{\textrm{ce}}(t) e^{j \aw_c t} } \} \notag \\
= & \calF \{ 0.5 \left( x_{\textrm{ce}}(t) e^{j \aw_c t} + x_{\textrm{ce}}^*(t) e^{-j \aw_c t} \right)  \} \notag \\
= & 0.5 \left( \calF \{ x_{\textrm{ce}}(t) e^{j \aw_c t} \} + \calF \{ x_{\textrm{ce}}^*(t) e^{-j \aw_c t} \} \right) \notag \\
= & 0.5 \left( X_{\textrm{ce}}(\aw - \aw_c) +  X_{\textrm{ce}}^*(-\aw - \aw_c) \right)
\label{eq:qam_complex_in_freq_domain}
\end{align}

\paragraph{Power of the complex envelope}
\[
P_{\textrm{ce}} = 2 P_{\textrm{QAM}}
\]

\section*{Analytic Signal and Hilbert Transform}

\paragraph{Inverse Fourier transform of the unit step}
\[
\calF^{-1}\{u(f)\} = \frac{\delta(t)}{2} + \frac{j}{2 \pi t} = \frac{1}{2} \left( \delta(t) + j \frac{1}{\pi t} \right)
\]

\paragraph{Analytic signal (frequency domain)}
\begin{equation}
X_{+}(f) = 2 u(f) X_{\textrm{QAM}}(f)
\label{eq:qamAnalytic}
\end{equation}

\paragraph{Analytic signal (time domain)}
\begin{align}
x_{+}(t) & = \calF^{-1} \{ X_{+}(f) \} = 2 \calF^{-1} \{u(f) X_{\textrm{QAM}}(f) \} \\
  & = x_{\textrm{QAM}}(t) + j \left( x_{\textrm{QAM}}(t) \conv \frac{1}{\pi t} \right)
\end{align}

\paragraph{Hilbert transform}
\[
\bar h(t) = \left\{ {\begin{array}{cc} {\frac 1 {\pi t}} & {, t \ne 0}\\ 0 & {, t=0} \end{array}  } \right.
\]
\[
\check{x}(t) = x(t) \conv \bar h(t) = \int_{- \infty}^\infty \frac{x(u)}{\pi (t - u)} du
\]
\[
x_{+}(t) = x_{\textrm{QAM}}(t) \conv \left(\delta(t) + j \frac 1 {\pi t} \right) = x_{\textrm{QAM}}(t) + j \check x(t)
\]

\section*{Sampling of Complex Signals}
\begin{equation}
\fs > F_\tmax
\label{eq:samplingTheoremComplexSignalsRepeated}
\end{equation}
(equivalently $\fs > \BW$ for a complex analytic baseband signal.)

\section*{Passband Channels and Baseband Equivalents}

\paragraph{Analytic passband impulse response and frequency response}
\begin{equation}
h_{+}(t) = h(t) + j \check{h}(t)
\label{eq:analyticPassbandTimeDomain}
\end{equation}
\begin{equation}
H_{+}(f) = 2 u(f) H(f)
\label{eq:analyticPassbandFrequencyDomain}
\end{equation}

\paragraph{Baseband equivalent channel}
\begin{equation}
h_{\textrm{bb}}(t) = h_{+}(t) e^{-j 2 \pi f_c t}
\label{eq:basebandEquivalentTimeDomain}
\end{equation}
\begin{equation}
H_{\textrm{bb}}(f) = H_{+}(f + f_c)
\label{eq:basebandEquivalentFreqDomain1}
\end{equation}
\begin{equation}
H_{\textrm{bb}}(f) = H(f + f_c)
\label{eq:eq:basebandEquivalentFreqDomain2}
\end{equation}

\paragraph{Baseband equivalent system}
\begin{equation}
y_{\textrm{bb}}(t) = x_{\textrm{bb}}(t) \conv \left( \frac{1}{2} h_{\textrm{bb}}(t) \right)
\end{equation}
\begin{equation}
Y_{\textrm{bb}}(f) = H(f+f_c) X_{\textrm{bb}}(f)
\end{equation}

\section*{Single-Tap Channel Estimation (Preamble-Based)}
\begin{equation}
\hat H(f_c) = \frac{1}{K} \sum_{k=1}^K \frac{Y_k}{X_k}
\end{equation}
(equalization is then done by normalizing the received symbols: $Y_k / \hat H(f_c)$.)
